%
%
%
%
%
%
%
%
%
%

%

%
%
%
\newcommand{\overbar}[1]{\mkern 1.5mu\overline{\mkern-1.5mu#1\mkern-1.5mu}\mkern 1.5mu}
%
%
\newcommand{\mybar}[3]{%
  \mathrlap{\hspace{#2}\overline{\scalebox{#1}[1]{\phantom{\ensuremath{#3}}}}}\ensuremath{#3}
}

\newcommand{\myhat}[3]{%
  \mathrlap{\hspace{#2}\widehat{\scalebox{#1}[1]{\phantom{\ensuremath{#3}}}}}\ensuremath{#3}
}

\newcommand{\barF}{\mybar{0.6}{2.5pt}{F}}
\newcommand{\barH}{\mybar{0.6}{2.5pt}{H}}
\newcommand{\barB}{\mybar{0.6}{2.5pt}{B}}
\newcommand{\barBs}{\mybar{0.6}{1.8pt}{B}}
\newcommand{\barPow}{\mybar{0.7}{1.5pt}{\Pow}}
\newcommand{\barSigmas}{\mybar{0.9}{0pt}{\Sigmas}}
\newcommand{\barrho}{\mybar{0.9}{1pt}{\rho}}

\newcommand{\gra}{\mathsf{gra}}
\newcommand{\dom}{\mathsf{dom}}

%
\DeclareMathOperator{\Sub}{\mathrm{Sub}}


\providecommand{\catname}{\mathsf} 
%
\providecommand{\clsname}{\mathcal}
\providecommand{\oname}[1]{{\mathop{\mathsf{#1}}\xspace}}

%
\def\defcatname#1{\expandafter\def\csname B#1\endcsname{\catname{#1}}}
\def\defcatnames#1{\ifx#1\defcatnames\else\defcatname#1\expandafter\defcatnames\fi}
\defcatnames ABCDEFGHIJKLMNOPQRSTUVWXYZ\defcatnames

%
\def\defclsname#1{\expandafter\def\csname C#1\endcsname{\clsname{#1}}}
\def\defclsnames#1{\ifx#1\defclsnames\else\defclsname#1\expandafter\defclsnames\fi}
\defclsnames ABCDEFGHIJKLMNOPQRSTUVWXYZ\defclsnames

%
%
%
%

%
\def\Set{\catname{Set}}
\def\Pos{\catname{Pos}}
\def\Nom{\catname{Nom}}
\def\Vect{\catname{Vec}_K}
\def\Conv{\catname{Conv}}
\def\Cpo{\catname{Cpo}}
\def\Cppo{\catname{Cppo}}
\def\dCpo{\catname{dCpo}}
\def\Top{\catname{Top}}

\newcommand{\Konig}{Kőnig}

%
\providecommand{\eps}{{\operatorname\epsilon}}
\providecommand{\amp}{\mathbin{\&}}
\providecommand{\from}{\leftarrow}

%
\providecommand{\argument}{-}
\providecommand{\altargument}{\underline{\;\;}}

\providecommand{\wave}[1]{\widetilde{#1}}		               %
\providecommand{\ul}{\underline}			                     %

\DeclareOldFontCommand{\bf}{\normalfont\bfseries}{\mathbf}
\providecommand{\mplus}{{\scriptscriptstyle\bf+}} 	       %

%
\providecommand{\PSet}{{\mathcal P}}			                 %
\providecommand{\FSet}{{\mathcal P}_{\omega}}		           %
\providecommand{\CSet}{{\mathcal P}_{\omega_1}}		         %
\providecommand{\NESet}{{\mathcal P}^{\mplus}}		         %
\providecommand{\Id}{\operatorname{Id}}

%
\providecommand{\Hom}{\mathsf{Hom}}
\providecommand{\id}{\mathsf{id}}
\providecommand{\op}{\mathsf{op}}
\providecommand{\comp}{\mathbin{\circ}}
\providecommand{\iso}{\mathbin{\cong}}
\providecommand{\tensor}{\mathbin{\otimes}}
\providecommand{\unit}{\star}				
\providecommand{\bang}{\operatorname!}				             %

%
\providecommand{\ito}{\hookrightarrow}				             %
\providecommand{\pto}{\mathrel{\rightharpoonup}}		       %
\providecommand{\tto}{\mathrel{\Rightarrow}}			         %
\providecommand{\mto}{\mapsto}
\providecommand{\xto}[1]{\,\xrightarrow{#1}\,}
\providecommand{\xfrom}[1]{\,\xleftarrow{\;#1}\,}
\providecommand{\To}{\mathrel{\Rightarrow}}			           %
\providecommand{\From}{\mathrel{\Leftarrow}}			         %
%
\providecommand{\flatxto}[1]{\xto{#1}}

%
%
%
%
%
%
%
%

\providecommand{\dar}{\kern-1.2pt\operatorname{\downarrow}}	
\providecommand{\uar}{\kern-1.2pt\operatorname{\uparrow}}	

%
\providecommand{\True}{\top}
\providecommand{\False}{\bot}

\providecommand{\bigor}{\bigvee}
\providecommand{\bigand}{\bigwedge}

\providecommand{\impl}{\Rightarrow}
\providecommand{\equ}{\Longleftrightarrow}

\providecommand{\entails}{\vdash}

%
\providecommand{\appr}{\sqsubseteq}
\providecommand{\join}{\sqcup}
\providecommand{\meet}{\sqcap}

\providecommand{\bigjoin}{\bigsqcup}
\providecommand{\bigmeet}{\bigsqcap}

%
\providecommand{\fst}{\oname{fst}}
\providecommand{\snd}{\oname{snd}}
\providecommand{\pr}{\prime}

\providecommand{\brks}[1]{\langle #1\rangle}
\providecommand{\Brks}[1]{\bigl\langle #1\bigr\rangle}

%
\providecommand{\inl}{\oname{inl}}
\providecommand{\inr}{\oname{inr}}
\providecommand{\inj}{\oname{in}}

\DeclareSymbolFont{Symbols}{OMS}{cmsy}{m}{n}
\DeclareMathSymbol{\iobj}{\mathord}{Symbols}{"3B}
%

%
\providecommand{\curry}{\oname{curry}}
\providecommand{\uncurry}{\oname{uncurry}}
\providecommand{\ev}{\oname{ev}}

%
\usepackage{stmaryrd}

\providecommand{\lsem}{\llbracket}
\providecommand{\rsem}{\rrbracket}
\providecommand{\sem}[1]{\lsem #1 \rsem}

\providecommand{\Lsem}{\bigl\llbracket}
\providecommand{\Rsem}{\bigr\rrbracket}
\providecommand{\Sem}[1]{\Lsem #1 \Rsem}

%
\providecommand{\comma}{,\operatorname{}\linebreak[1]}		 %
\providecommand{\dash}{\nobreakdash-\hspace{0pt}}		       %
\providecommand{\erule}{\rule{0pt}{0pt}}		               %
                                                           %

\providecommand{\by}[1]{\text{/\!\!/~#1}}			             %
\providecommand{\pacman}[1]{}					                     %

\newcommand{\undefine}[1]{\let #1\relax}					                       %

\providecommand{\noqed}{\def\qed{}}				                 %

%
\providecommand{\mone}{{\text{\kern.5pt\rmfamily-}\mathsf{\kern-.5pt1}}}

\providecommand{\smin}{\smallsetminus}

%
%
%
%
%

%
%
%
%

%
%
%

%
%
 
%
\makeatletter
\def\mfix#1{\oname{#1}\@ifnextchar\bgroup\@mfix{}}	       %
\def\@mfix#1{#1\@ifnextchar\bgroup\mfix{}}			           %
\makeatother

%
\providecommand{\ift}[3]{\mfix{if}{\mathbin{}#1}{then}{\mathbin{}#2}{else}{\mathbin{}#3}}
\providecommand{\case}[3]{\mfix{case}{\mathbin{}#1}{of}{#2}{\kern-1pt;}{\mathbin{}#3}}




\usepackage{dsfont}
%

\DeclareMathSymbol{\mathinvertedexclamationmark}{\mathord}{operators}{'074}
\DeclareMathSymbol{\mathexclamationmark}{\mathord}{operators}{'041}
\makeatletter
\newcommand{\raisedmathinvertedexclamationmark}{%
  \mathord{\mathpalette\raised@mathinvertedexclamationmark\relax}%
}
\newcommand{\raised@mathinvertedexclamationmark}[2]{%
  \raisebox{\depth}{$\m@th#1\mathinvertedexclamationmark$}%
}
\makeatother

\newcommand{\sqdown}{\logp{\Downarrow}}

\newcommand{\lfplong}{locally finitely presentable\xspace}
\newcommand{\Lfplong}{Locally finitely presentable\xspace}
\newcommand{\lfp}{\lfplong}
\newcommand{\Lfp}{\Lfplong}

\newcommand{\strength}{\mathsf{st}}
\newcommand{\product}{\times}
\newcommand{\wt}{\widetilde}
\newcommand{\longTo}{\Longrightarrow}
\newcommand{\Pt}{V}
\newcommand{\R}{\mathcal{R}}
%
\newcommand{\lsmc}{\fatsemi}
\newcommand{\Pre}{\mathbf{Pre}}
\newcommand{\Pow}{\mathcal{P}}
\newcommand{\xTo}{\xRightarrow}
\newcommand{\HJBis}{\mathbf{HJBisim}}
\newcommand{\Bis}{\mathbf{Bisim}}
\newcommand{\primrec}{\mathsf{pr}\,}
\newcommand{\hattau}{\hat{\tau}}
\newcommand{\hatini}{\hat{\ini}}
\newcommand{\hatr}{\hat{r}}
\newcommand{\hatbeta}{\hat{\beta}}
\newcommand{\lft}{{\mathsf{l}}}
\newcommand{\rgt}{{\mathsf{r}}}
\DeclareMathOperator{\colim}{\mathsf{colim}}
\newcommand{\terminal}{\mathexclamationmark}
\newcommand{\initial}{\raisedmathinvertedexclamationmark}
\newcommand{\smc}{\mathbin{;}}
\newcommand{\hatR}{{\widehat{R}}}
\newcommand{\qand}{\quad\text{and}\quad}
\newcommand{\qqand}{\qquad\text{and}\qquad}
\newcommand{\rsto}{\rightsquigarrow}
\newcommand{\pb}[2]{\scalebox{#1}{\rotatebox{#2}{$\ulcorner$}}}
\newcommand{\deq}[1]{{\rotatebox{#1}{$=$}}}
\newcommand{\dleq}[1]{{\rotatebox{#1}{$\preceq$}}}
\newcommand{\dgeq}[1]{{\rotatebox{#1}{$\succeq$}}}
\newcommand{\cev}[1]{\reflectbox{\ensuremath{\vec{\reflectbox{\ensuremath{#1}}}}}}
\newcommand{\diag}{\mathsf{in}}
\newcommand{\HO}{\mathcal{HO}}
\newcommand{\Refl}{\mathbf{Refl}}
\newcommand{\Gra}{\mathbf{Gra}}
\newcommand{\RGra}{\mathbf{RGra}}
\newcommand{\RRel}{\mathbf{RRel}}
\newcommand{\happrox}{\mathrel{\widehat\lesssim}}
\newcommand{\happroxs}{\mathrel{\hat\approx^{\star}}}
\newcommand{\under}[1]{\lvert#1\rvert}
\newcommand{\SKI}{\mathrm{SKI}}
\renewcommand{\L}{\mathcal{L}}
\newcommand{\ap}{{\mathrm{ap}}}
\newcommand{\var}{\mathsf{var}}
\newcommand{\cn}{{\mathrm{cn}}}
\newcommand{\cv}{{\mathrm{cv}}}
%
%
\newcommand{\Powf}{\mathcal{P}_{\omega}}
\newcommand{\FSigma}{{\mathbb{\Sigma}}}
\newcommand{\FSigmas}{{\mathbb{\Sigma}}^\star}
\newcommand{\Sigmas}{\Sigma^{\star}}
\newcommand{\FDelta}{{\mathbb{\Delta}}}
\newcommand{\FDeltas}{{\mathbb{\Delta}}^\star}
\newcommand{\ar}{\mathsf{ar}}
%
\newcommand{\NT}{\mathrm{Nat}}
\newcommand{\supp}{\mathsf{supp}}
\newcommand{\epito}{\twoheadrightarrow}
\newcommand{\abs}{\mathsf{abs}}
\newcommand{\conc}{\mathsf{conc}}
\newcommand{\T}{{\mathcal{T}}}
\renewcommand{\S}{{\mathcal{S}}}
\newcommand{\can}{\mathsf{can}}
\newcommand{\lab}{\mathsf{label}}
\newcommand{\leaves}{\mathsf{leaves}}
\newcommand{\incl}{\mathsf{in}}
\newcommand{\seq}{\subseteq}
\newcommand{\ol}{\overline}
\newcommand{\out}{\mathsf{out}}
\newcommand{\outl}{\mathsf{l}}
\newcommand{\outr}{\mathsf{r}}
\newcommand{\trc}{{\mathsf{trc}}}
\newcommand{\fn}{{\mathsf{fn}}}
\newcommand{\beh}{{\mathsf{beh}}}
\newcommand{\dbtilde}[1]{\widetilde{\raisebox{0pt}[0.85\height]{$\widetilde{#1}$}}}
\providecommand{\C}{}
\providecommand{\D}{}
\providecommand{\E}{}

\newcommand{\fp}{{\mathrm{fp}}}
\newcommand{\fg}{{\mathrm{fg}}}
\newcommand{\rec}{{\mathrm{rec}}}
\newcommand{\wf}{{\mathrm{wf}}}
\newcommand{\frc}{\Coalg_{\fp,\rec}(H)}
\newcommand{\fwf}{\Coalg_{\fp,\wf}(H)}
\newcommand{\fgwf}{\Coalg_{\fp,\wf}(H)}

%
%
%

\renewcommand{\C}{{\mathbb{C}}}
\renewcommand{\D}{{\mathbb{D}}}
\renewcommand{\E}{{\mathbb{E}}}
\newcommand{\I}{{\mathbb{I}}}

\newcommand{\A}{\mathcal{A}}

\newcommand{\iref}[2]{\autoref{#1}.\ref{#1:#2}}


\renewcommand{\id}{{\mathsf{id}}}
\newcommand{\Nat}{\mathds{N}}
\renewcommand{\NT}{\mathsf{Nat}}
\newcommand{\Rel}{\mathbf{Rel}}

\newcommand{\ngoesv}[2]{\ensuremath{#1 \centernot{\xrightarrow{~#2~}}}} %

\newcommand{\Goes}[2]{\ensuremath{#1 \Rightarrow #2}}
\newcommand{\goes}[2]{\ensuremath{#1 \rightarrow #2}}
%
%
\newcommand{\Rets}[2]{\ensuremath{#1 \Downarrow #2}}
\newcommand{\Retsc}[2]{\ensuremath{#1 \Rightarrow #2, \checkmark}}
 \newcommand{\rets}[2]{\ensuremath{#1 \downarrow #2}}
\newcommand{\retsc}[2]{\ensuremath{#1 \rightarrow #2, \checkmark}}
%
%

%
\newcommand{\Goesv}[3]{\ensuremath{#1 \xRightarrow{~#3~} #2}}
\newcommand{\goesv}[3]{\ensuremath{#1 \xrightarrow{~#3~} #2}}
%
%
%
\newcommand{\retsv}[3]{\ensuremath{#1 \xrightarrow{~#3~} #2, \checkmark}}
%
%

\newcommand{\retsvr}[3]{\ensuremath{#1~\textcolor{red}{\Downarrow_{#3}}~#2}}
%
%


%
%
\newcommand{\stimes}[1]{{S\product #1}}
\newcommand{\xtimes}[1]{{X\product #1}}
\newcommand{\ytimes}[1]{{Y\product #1}}
\newcommand{\sbrks}[1]{{ \brks{\fst,#1}  }}
\newcommand{\cpair}[2]{{ #1, #2 }}
\newcommand{\tpair}[2]{{ #1 \product #2 }}


%
%
%
\renewcommand{\brack}[1]{\lbrack #1 \rbrack}
%
\newcommand{\sbrack}[1]{\llbracket #1 \rrbracket}
%
\newcommand{\ssbrack}[1]{\llbracket\mkern-4.7mu\llbracket #1 \rrbracket\mkern-4.7mu\rrbracket}
%
\newcommand{\f}{\oname{f}}
\newcommand{\fii}{\oname{f}}
%

%
\newcommand{\takeout}[1]{\empty}

\newcommand{\Dist}{\mathcal{D}}
\newcommand{\CPowf}{\Pow_{\mathsf{c},\omega}}

\newcommand{\ini}{\iota}
\newcommand{\ter}{\tau}
\newcommand{\wh}{\widehat}
\DeclareMathOperator{\Coalg}{\mathsf{Coalg}}
\DeclareMathOperator{\Alg}{\mathsf{Alg}}
%
\newcommand{\cofun}{(\overline{\rule{-1pt}{7pt}\argument})}
\renewcommand{\rho}{\varrho}
\newcommand{\N}{\mathds{N}}
\newcommand{\opp}{\mathsf{op}}
\newcommand{\xla}[1]{\xleftarrow{~#1~}}
\newcommand{\mypowfin}{\mathscr{P}_\omega}

%
\newcommand{\pullbackangle}[2][]{\arrow[phantom,to path={
                     -- ($ (\tikztostart)!1cm!#2:([xshift=8cm]\tikztostart) $)
                        node[anchor=west,pos=0.0,rotate=#2,
                        inner xsep = 0]
                        {\begin{tikzpicture}[minimum
                        height=1mm,baseline=0,#1]
    \draw[-] (0,0) -- (.5em,.5em) -- (0,1em);
                        \end{tikzpicture}}}]{}}

\makeatletter
\newsavebox{\@brx}
\newcommand{\llangle}[1][]{\savebox{\@brx}{\(\m@th{#1\langle}\)}%
  \mathopen{\copy\@brx\kern-0.5\wd\@brx\usebox{\@brx}}}
\newcommand{\rrangle}[1][]{\savebox{\@brx}{\(\m@th{#1\rangle}\)}%
  \mathclose{\copy\@brx\kern-0.5\wd\@brx\usebox{\@brx}}}
\makeatother

\renewcommand{\comp}{\cdot}
\renewcommand{\c}{\colon}

\newcommand{\Var}{\mathop{\mathsf{Var}}}
\newcommand{\xra}{\xrightarrow}
\renewcommand{\xto}{\xra}
\let\xmpsto=\xmapsto
%

\newcommand{\V}{\mathcal{V}}
\newcommand{\monoto}{\rightarrowtail}
\newcommand{\subto}{\hookrightarrow}
\newcommand{\commu}{\ensuremath{\circlearrowleft}}
\renewcommand{\Nat}{\mathbb{N}}
\newcommand{\pfun}{\rightharpoonup}
\newcommand{\fset}{{\mathbb{F}}}
\newcommand{\vcat}{{\Set}^{\fset}}
\newcommand{\alg}[1]{\mathbf{Alg}(#1)}
\newcommand{\coalg}[1]{\mathbf{CoAlg}(#1)}
\newcommand{\mon}{\bullet}

\newcommand{\act}[1]{\mathbf{Act}(#1)}
\newcommand{\adj}[4]{#1\negmedspace: #2\rightleftarrows #3:\negmedspace #4}
\newcommand{\gcat}{\mathbb{C}}
\providecommand{\Dinat}{\mathsf{Dinat}}
\providecommand{\DiNat}{\mathsf{DiNat}}
\newcommand{\mS}{{\mu\Sigma}}
\DeclareMathOperator{\iter}{\oname{it}}
\DeclareMathOperator{\coiter}{\oname{coit}}
\newcommand{\transpose}[1]{\overline{#1}}
\renewcommand{\epsilon}{\varepsilon}

\newcommand*\xbar[1]{%
  \kern.2em\hbox{%
    \vbox{%
      \hrule height 0.5pt %
      \kern0.5ex%
      \hbox{%
        \kern-0.2em%
        \ensuremath{#1}%
        \kern-0.4em%
      }%
    }%
  }\kern.4em %
} 

\makeatletter
\newcommand{\monto}{\@ifstar{\@mtolifted}{\@mto}}
\newcommand{\@mto}{\multimapdot}
\newcommand{\@mtolifted}{\mathbin{\xbar{\multimapdot}}}
\makeatother

\newcommand{\Ar}[1]{\mathbf{Ar}(#1)}

\newcommand{\bool}{\mathsf{bool}}
\newcommand{\true}{\mathsf{true}}
\newcommand{\false}{\mathsf{false}}
\newcommand{\ite}[3]{\mathsf{if}\,{#1}\,\mathsf{then}\,#2\,\mathsf{else}\,#3}
\newcommand{\app}[2]{\,}
\newcommand{\skiterm}{\Lambda_{\varsigma}}
\newcommand{\skitermu}{\Lambda_{u}}
\newcommand{\finalc}{Z}

%
%
%
%
